\documentclass[a4paper,twoside]{article}
\usepackage{graphicx} % Required for inserting images
\usepackage[hidelinks]{hyperref} % geen kaders om linken
\usepackage{parskip}
\usepackage{bm} % For bolt symbols eg \bm{\nabla}
\usepackage{amsmath, amssymb, amsfonts} % For nicer math format
\usepackage{esint} % For bolletje in integrals
\usepackage{mathrsfs} % For curly arrr
\usepackage{subcaption}
\usepackage[english]{babel}
\usepackage{dsfont}
\usepackage[a4paper, margin=0.24cm, landscape]{geometry} 
% Verander 1.5cm naar bijvoorbeeld 1cm voor een extreem vol blad.
\usepackage[T1]{fontenc}
\usepackage[utf8]{inputenc}
\usepackage{lmodern}
\usepackage{multicol}

%Tikz packages
\usepackage{tikz}
\usetikzlibrary{angles,quotes, babel} % for pic (angle labels)
\usetikzlibrary{arrows.meta,decorations.markings,calc, decorations.pathmorphing, decorations.pathreplacing, positioning, shapes, shadings, patterns}
\usepackage{xcolor}
\usepackage{tikz-3dplot}
\usepackage[siunitx]{circuitikz}


\usepackage{etoolbox} % ifthen
\usepackage[outline]{contour} % glow around text
\usepackage{xfp} % higher precision (16 digits?)
\contourlength{1.1pt}

\setlength{\parindent}{0pt}
\setlength{\parskip}{0.08ex}
\setlength{\columnsep}{0.12cm}
%\linespread{0.8} 
%\pagestyle{empty}

\newcommand{\secblock}[1]{\vspace{0.08ex}\textbf{#1}\par}
\newcommand{\eqline}[1]{\(#1\)\par}


%\usepackage{geometry} %size regulating
%\geometry{
 %a4paper,
 %left=20mm,
 %right=20mm,
 %top=10mm,
 %bottom=20mm,
%}

% Source - https://tex.stackexchange.com/a/506180
% Posted by Henri Menke, modified by community. See post 'Timeline' for change history
% Retrieved 2026-05-07, License - CC BY-SA 4.0

%vptovf griffm.vpl
%vptovf griffb.vpl

%\DeclareFontFamily{U}{griff}{}
%\DeclareFontShape{U}{griff}{m}{n}{<-> s*[2.2] griffm}{}
%\DeclareFontShape{U}{griff}{b}{n}{<-> s*[2.2] griffb}{}
%\DeclareSymbolFont{griff}{U}{griff}{m}{n}
%\SetSymbolFont{griff}{bold}{U}{griff}{b}{n}
%\DeclareMathSymbol{\rcurs}{\mathalpha}{griff}{"72}
%\DeclareBoldMathCommand{\brcurs}{\rcurs}
%\newcommand*\hrcurs{\hat{\brcurs}}


\begin{document}


%\fontsize{4.7}{4.95}\selectfont Om alles heel kleini te maken

\begin{multicols*}{4}
\secblock{H1 Vector Analysis}
%Meeste staat op formularium
\secblock{H2 Electrostatics}

Coulomb Force: 
\(\bm F = \frac{1}{4\pi \epsilon_0}\frac{qQ}{r^2}\hat{r}\)

Electric field:
\(
\bm E (\bm r) = \frac{1}{4\pi \epsilon _0} \int \frac{\rho(\bm r ')}{r^2}\hat{\bm r} d\tau'
\)

Electric flux: 
\(
\Phi_E = \int_S \bm E \cdot d\bm a
\)

Closed path:
\(
\oint \bm E \cdot d \bm l = 0 \implies \bm\nabla
 \times \bm E = \bm 0
\)

Electric potential:
\(V(\bm r) =-\int_0^r \bm E \cdot d \bm l \implies \bm E = - \nabla V \)

% Poison/Laplace equation???

Work: 
\(W = \int _a^b \bm F \cdot d \bm l = -Q \int _a^b \bm E \cdot d \bm l  \implies W = QV(\bm r)\)

or:
\[
W = \frac{\epsilon_0}{2}\int E^2d\tau
\]

%Electrostatic pressure ? (2.51)

Capacitance:
\[
C = \frac{Q}{V} \implies W = \frac{1}{2}CV^2
\]



\secblock{H3 Potentials}

%Laplace again (maar 3.1 niet te kennen?) (3.1)

% Method of imagees (3.2) 

% Separation oof variables (3.3)

% Afleiding Legendre polynomials?

%Rodrigues formula:
%\[
%\bm P_l(x) = \frac{1}{2^l l! }(\frac{d}{dx})^l (x^2-1)^l
%\]

Curly r:

\(\frac{1}{r_{curly}} = \frac{1}{r}\sum^\infty_{n=0} (\frac{r'}{r})^n P_n (\cos \alpha)\)

Dipole contribution to potential:

\(V_{\text{dip}} (\bm r) = \frac{1}{4 \pi \epsilon_0} \frac{\bm p \cdot \hat{\bm r}}{r^2}\)

\(\bm E_{dip}(r,\theta)=\frac{p}{4\pi \epsilon_0 r^3}(2\cos \theta \hat{\bm r} + \sin \theta \hat{\bm \theta})\)



\secblock{H4 Electric Fields in Matter}

Atomic polarizability:
\[
\bm p = \alpha\bm E
\]

Torque:
\[
\bm N = q \bm d \times \bm E = \bm p \times \bm E
\]
\[
\bm F = (\bm p \cdot \bm \nabla)\bm E
\]
Bounded surface charge density:
\[ \bm P \cdot \hat{\bm n}\]
Bounded volume charge density:
\[
\rho_b = - \bm \nabla \cdot \bm P
\]
Potential Macroscopic field:
\[
V (\bm r) = \frac{1}{4 \pi \epsilon_0} \int \frac{\bm P(\bm r') \cdot \hat{\bm r_{c}}}{r^2_c} d\tau'
\]

Electric displacement integral form:
\[
\oint \bm D \cdot d\bm a = Q_{f_{enc}}
\]

Electric field in a homogeneous linear dielectric media:
\[
\bm E = \frac{1}{4\pi \epsilon _0} \frac{q}{r^2}\hat{\bm r}
\]
Bound charge in homogeneous isotropic linear dielectric:
\[
\rho _b = - \frac{\chi_e}{1+\chi_e} \rho_f
\]

Work:
\[
\frac{1}{2} \int \bm D \cdot \bm E d\tau
\]

% Capacitance????






\secblock{H5 Magnetostatics}
Magnetic forces do no work!\\
\(
\bm F_{mag} = I \int d \bm l \times \bm B
\)

Surface current:

\(
\bm K = \frac{d \bm I}{dl_\perp} = \sigma \bm v
\)

Magnetic force on the surface current

\(
\bm F_{mag} =  \int ( \bm v \times \bm B)\sigma da =  \int ( \bm K \times \bm B)da
\)

Volume current density:

\(
\bm J = \frac{d \bm I}{dl_\perp} = \rho \bm v
\)

Magnetic force on the volume current:

\(
\bm F_{mag} =  \int ( \bm v \times \bm B)\rho d\tau =  \int ( \bm J \times \bm B)d\tau
\)

Continuity equation:

\(\bm \nabla \cdot \bm J = -\frac{\partial \rho}{\partial t}\)

Biot-Savart law:

\(
\bm{B}(\bm{r})=\frac{\mu_0}{4\pi}\int\frac{I\times\hat{\bm{r}}_c}{r_c^2}d\bm l'
=\frac{\mu_0 I}{4\pi}\int\frac{dl'\times\hat{\bm{r}}_c}{r^2_c}
\)

For surface currents:

\(\bm{B}(\bm{r})=\frac{\mu_0}{4\pi}\int\frac{\bm K(\bm r')\times\hat{\bm{r}}_c}{r_c^2}d a'\)

For volume currents:

\(\bm{B}(\bm{r})=\frac{\mu_0}{4\pi}\int\frac{\bm J(\bm r')\times\hat{\bm{r}}_c}{r_c^2}d \tau '\)

Vector potential in magnetostatics:

\(\bm A = \frac{\mu_0}{4\pi}\int\frac{\bm J(\bm r')}{r_c}d \tau '\).

\(A_{dip}(\bm r) = \frac{\mu_0}{4\pi}\frac{\bm m \times \hat{\bm r}}{r^2}\)

magnetic dipole moment:

\(\bm m = I \bm a\)

\(\bm A_{dip} (\bm r) = \frac{\mu_0}{4\pi} \frac{m \sin \theta}{r^2}\hat{\bm \phi}\)

\(\bm B_{dip} (\bm r) = \bm \nabla \times \bm A_{dip} (\bm r) = \frac{\mu_0 m}{4\pi r^3}(2\cos \theta \hat{\bm r} + \sin \theta \hat{\bm \theta})\)
\secblock{H6 Magnetic Fields in Matter}
Torque:
\(\bm N = \bm m \times \bm B\)
Force nonuniform field "fringing":
\(\bm F = 2\pi \bm I \bm R \bm B \cos\theta\)
Force nonuniform field: 
\(\bm F = \bm \nabla (\bm m \cdot \bm B)\)

% Effect of a Magnetic field on atomic orbits?

Vector potential of a single magnetic dipole:

\(\bm A (\bm r) = \frac{\mu_0}{4\pi} \frac{\bm m \times \hat{\bm r}_c}{r^2_c} = \frac{\mu_0}{4\pi} \int \frac{\bm M (\bm r) \times \hat{\bm r}_c}{r^2_c} d\tau'\)

The potential of a volume current:

\(J_b = \bm \nabla \times \bm M\)

The potential of a surface current:

\(\bm K_b = \bm M \times \hat{\bm n}\)  



\secblock{H7 Electrodynamics}

Current density (Ohm's law):

\(\bm J = \sigma \bm f = \sigma (\bm E + \bm v \times \bm B) \approx \sigma \bm E\)

Resistivity:

\(\rho = \frac{1}{\sigma}\)

Perfect conductor:

\(\sigma = \infty\) or/and \(\bm E = \frac{\bm J}{\sigma}=0\)

Perfect isolator:
\(\sigma = 0\)

\(V = IR\)

For steady currents and uniform conductivity:

\(\bm \nabla \cdot \bm E = \frac{1}{\sigma} \bm \nabla \cdot \bm J = 0\)

% Rare formule (7.6)??

Joule heating law:
\(P = VI = I^2R\)

Electromotive force (emf):

\(\epsilon_{emf} = \oint\bm f \cdot d \bm l\)

Flux rule:
\(\epsilon_{emf} = -\frac{d \phi}{dt}\)

Analog to Biot-Savart law:

\(\bm{E}(\bm{r})=-\frac{1}{4\pi}\frac{\partial}{\partial t}\int\frac{\bm B\times\hat{\bm{r}}_c}{r_c^2}d\bm \tau'\)

Inductance:

\(\bm B_1 = \frac{\mu_0 I_1}{4\pi} \oint \frac{dl_1 \times \hat{\bm{r}}_c}{r_c^2}\)

Mutual inductance:

\(\Phi_{21} = \int \bm B_1 \cdot d\bm a_2 = M_{21} I_1\)

Neuman's law:
\(M_{21} = M_{12}\)

\(M_{21} = \frac{\mu_0}{4\pi} \oint \oint \frac{d\bm l_1 \cdot d\bm l_2}{r_c}\)

Self inductance:

\(\Phi_{11} = \int \bm B_1 \cdot d\bm a_1 = L I_1 \implies \Phi = LI\)

emf of self inductance:

\(\epsilon_{emf} = -L \frac{dI}{dt}\)

Power:

\(\frac{dW}{dt} = -\epsilon_{emf} I = LI \frac{dI}{dt}\)

Work done:
\(W = \frac{1}{2} L I^2\)

\(W = \frac{1}{2\mu_0} \int_{all \, space} B^2 d\tau\)

\(W_{elec} = \frac{1}{2} \int  (V \rho) d\tau = \frac{\epsilon_0}{2} \int  E^2 d\tau\)

\(W_{mag} = \frac{1}{2} \int (\bm A \cdot \bm J) d\tau = \frac{1}{2\mu_0} \int B^2 d\tau\)

Displacement current (in medium):

\(\bm J_d = \epsilon_0 \frac{\partial \bm D}{\partial t}\)

Polarization current (in medium):

\(\bm J_p = \frac{\partial \bm P}{\partial t}\)



\secblock{H8 Conservation Laws}

Local charge conservation:

\(\bm \nabla \cdot \bm J + \frac{\partial \rho}{\partial t} = 0\)

%Extra
%momentum density:
%\(\bm g = \epsilon_0 \bm E \times \bm B\)

Total energy stored in the electromagnetic field:

\(u = \frac{1}{2} (\epsilon_0 E^2 + \frac{1}{\mu_0} B^2)\)

Work done on charge q:

\(dW = q \bm E \cdot \bm v dt\)

Poynting theorem:

\(\frac{\partial u}{\partial t} + \bm \nabla \cdot \bm S = -\bm J \cdot \bm E\)

Poynting vector:

\(\bm S = \frac{1}{\mu_0} \bm E \times \bm B\)

Continuity equation for energy:

\(\frac{\partial u}{\partial t} + \bm \nabla \cdot \bm S = 0\)


%Afleiding momentum density en Maxwell stress tensor H8.2.2 nog inzetten??

Maxwell stress tensor:

\(T_{ij} = \epsilon_0 (E_iE_j-\delta_{ij}E^2) +\frac{1}{\mu_0} (B_iB_j-\delta_{ij}B^2)\)

Force per unit volume:

\(\bm f = \bm\nabla \cdot \overleftrightarrow{T} - \epsilon_0 \mu_0 \frac{\partial \bm S}{\partial t}\)

Total electromagnetic force on a charge distribution:

\(\bm F = \int \bm f d\tau = \oint_S \overleftrightarrow{T} \cdot d\bm a - \epsilon_0 \mu_0 \frac{d}{dt} \int_V \bm S d\tau\)   

At the boundary:

\(\bm F = \oint _S \overleftrightarrow{T} \cdot d\bm a\) (static case)

%Niet te kennen deel 8.2.3 en verder??

%Momentum stored in the EM field:
%\(\bm p = \frac{1}{c^2} \int \bm S d\tau = \mu_0 \epsilon_0 \int \bm S d\tau  \)

%Conservation of momentum:
%\(\frac{d\bm p}{dt} + \oint_S \overleftrightarrow{T} \cdot d\bm a = 0\)

%Momentum density:   
%\(\bm g = \frac{1}{c^2} \bm S = \mu_0 \epsilon_0 \bm S\)

%Continuity equation for momentum:
%\(\frac{\partial \bm g}{\partial t} + \bm \nabla \cdot \overleftrightarrow{T} = 0\)


\secblock{H9 Electromagnetic Waves}

Energy per unit volume:

\(u = \frac{1}{2} (\epsilon_0 E^2 + \frac{1}{\mu_0} B^2)\)

Case monochromatic plane wave:

\(B^2 = \frac{1}{c^2} E^2 = \epsilon_0 \mu_0 E^2\)

Electric and magnetic contributions to energy density are equal:

\(u = \epsilon_0 E^2 = \epsilon_0 E_0^2 \cos^2(k z - \omega t + \delta)\)

Poynting vector becomes:

\(\bm S = c \epsilon_0 E^2 \cos^2(k z - \omega t + \delta) \hat{\bm z} = c u \hat{\bm z}\)

Momentum density stored in field:

\(\bm g = \frac{1}{c^2} \bm S = \frac{1}{c} u \hat{\bm z}\)

Averages: 

\(\langle u \rangle = \frac{1}{2} \epsilon_0 E_0^2\)

\(\langle \bm S \rangle = \frac{1}{2} c \epsilon_0 E_0^2 \hat{\bm z}\)

\(\langle \bm g \rangle = \frac{1}{2} \frac{u}{c} \hat{\bm z}\)

Intensity: \(\bm I \equiv \langle  S \rangle = \frac{1}{2} c \epsilon_0 E_0^2 \)

Radiation pressure:

\(P = \frac{1}{A}\frac{\Delta p}{\Delta t} = \frac{1}{2}\epsilon_0 E_0^2 =\frac{I}{c}\)

Speed of light in medium:
\(v = \frac{c}{n}\)

Where: \(n \equiv \sqrt{\frac{\epsilon \mu}{\epsilon_0 \mu_0}} \approx \sqrt{\epsilon_r}\)

The energy density in a medium:

\(u = \frac{1}{2} (\epsilon E^2 + \frac{1}{\mu} B^2)\)

Intensity in a medium:
\(I = \frac{1}{2} v \epsilon E_0^2 \)

boundary conditions:

\(\epsilon_1 E_1^{\perp} = \epsilon_2 E_2^{\perp}\)

\(\frac{1}{\mu_1} B_1^{\perp} = \frac{1}{\mu_2} B_2^{\perp}\)

\secblock{H10 Potentials and Fields}
\(\bm B\) is divergenceless, so:
\(\bm B = \bm \nabla \times \bm A\)

and: \(\bm E = - \bm \nabla V - \frac{\partial \bm A}{\partial t}\)

\(\nabla^2 V + \frac{\partial }{\partial t}(\bm \nabla \cdot \bm A) = -\frac{\rho}{\epsilon_0}\)

\(\nabla \times (\nabla \times \bm A)= \mu_0 \bm J - \mu_0 \epsilon_0 \bm \nabla \left(\frac{\partial V }{\partial t}\right) - \mu_0 \epsilon_0 \frac{\partial^2 \bm A}{\partial t^2}\)

\( \left( \nabla^2 \bm A - \mu_0 \epsilon_0 \frac{\partial^2 \bm A}{\partial t^2} \right) 
- \bm \nabla \left(\bm \nabla \cdot \bm A + \mu_0 \epsilon_0 \frac{\partial V}{\partial t} \right)
= -\mu_0 \bm J\)

Gauge transformation:
\(\bm A' = \bm A + \bm \alpha\) and \(V' = V + \beta\) 

\(\bm \nabla \times \bm\alpha = 0\)

\(\bm \alpha = \bm \nabla \lambda\)

\(\bm \nabla \beta + \frac{\partial \bm \alpha}{\partial t} = 0\)

\(\bm \nabla (\beta + \frac{\partial \lambda}{\partial t})=0\)

\(\bm A' = \bm A + \bm \nabla \lambda\) and \(V' = V - \frac{\partial \lambda}{\partial t}\)

For any function \(\lambda(\bm r, t)\).

Coulomb gauge: \(\bm \nabla \cdot \bm A = 0\)

Lorenz gauge: \(\bm \nabla \cdot \bm A = -\mu_0 \epsilon_0 \frac{\partial V}{\partial t} \)

then: \(\nabla^2 -\mu_0 \epsilon_0 \frac{\partial^2}{\partial t^2} \equiv \Box^2\)

d'Alembertian operator: 

\(\Box^2 V = -\frac{\rho}{\epsilon_0}\) and \(\Box^2 \bm A = -\mu_0 \bm J\)

Lorentz force in terms of potentials:

\(\bm F = q \left[ -\bm \nabla V - \frac{\partial \bm A}{\partial t} + \bm v \times (\bm \nabla \times \bm A) \right]\)
\(= -q[\frac{\partial \bm A}{\partial t} + (\bm v \cdot \bm \nabla) \bm A + \bm \nabla (V-\bm v \cdot \bm A)]\)

gives us:

\( \frac{d \bm A}{dt} = \frac{\partial \bm A}{\partial t} + (\bm v \cdot \bm \nabla) \bm A \)

\(\frac{d\bm p}{dt} = - \bm \nabla U\)

Canonical momentum:
\(\bm p_{can} = \bm p + q \bm A\)

\(U_{vel} = q(V-\bm v \cdot \bm A)\)

Static case:
\(\nabla^2 V = -\frac{\rho}{\epsilon_0}\) and \(\nabla^2 \bm A = -\mu_0 \bm J\)

retarded time:
\(t_r = t - \frac{r_c}{c}\)

retarded potentials:

\(V(\bm r, t) = \frac{1}{4 \pi \epsilon_0} \int \frac{\rho(\bm r', t_r)}{r_c} d\tau'\)
and

\(\bm A(\bm r, t) = \frac{\mu_0}{4\pi} \int \frac{\bm J(\bm r', t_r)}{r_c} d\tau'\)








\secblock{H11 Radiation}

\secblock{H12 Electrodynamics and Relativity}

\secblock{Handige rekenregels}

Coordinate transoformation:

Spherical:

\(\hat{\bm x} = \sin \theta \cos \phi \hat{\bm r} + \cos \theta \cos \phi \hat{\bm \theta}-\sin \phi \hat{\bm \phi}\)

\(\hat{\bm y} = \sin\theta \sin \phi \hat{\bm r} + \cos\theta \sin \phi \hat{\bm \theta} + \cos \phi \hat{\bm \phi}\)

\(\hat{\bm z} = \cos \theta \hat{\bm r} - \sin \theta \hat{\bm \theta}\)

\(\hat{\bm{r}} = \sin\theta \cos \phi \hat{\bm x} + \sin \theta \sin \phi \hat{\bm y}+\cos \theta \hat{\bm z}\)

\(\hat{\bm \theta} = \cos\theta\cos\phi \hat{\bm x} + \cos\theta\sin\phi \hat{\bm y} - \sin\theta \hat{\bm z}\)

\(\hat{\bm \phi} = -\sin\phi \hat{\bm x} + \cos\phi \hat{\bm y}\)

Cylindrical:

\(\hat{\bm x} = \cos \phi \hat{\bm s} - \sin \phi \hat{\bm \phi}\)

\(\hat{\bm y} = \sin \phi \hat{\bm s} + \cos \phi \hat{\bm \phi}\)

\(\hat{\bm z} = \hat{\bm z}\)

\(\hat{\bm s} = \cos \phi \hat{\bm x} + \sin \phi \hat{\bm y}\)

\(\hat{\bm \phi} = -\sin \phi \hat{\bm x} + \cos \phi \hat{\bm y}\)

Goniometric identities:





\end{multicols*}
\end{document}
